\subsection{Regular rings and finite algebras}

PROPOSITION 5.-- Let \(\rho: A \to B\) be a homomorphism of Noetherian rings and N a B-module. Suppose that

\begin{enumerate}
\def\labelenumi{\alph{enumi})}
\item
  the ring A is regular,
\item
  N is a finitely generated A-module,
\item
  the B-module N is Macaulay,
\item
  every minimal prime ideal of Supp \(_{B}\) (N) lies over a minimal prime ideal of A.
\end{enumerate}

Then N is a projective (finitely generated) A-module.

It is a question of proving that, for every maximal ideal m of A, the \(A_{m}\) -module \(N_{m}\) is free (II, § 5, no. 2, th. 1). The A-module \(B/Ann_{B}(N)\) is a submodule of the finitely generated A-module \(\operatorname{End}_{A}(N)\) , hence is finitely generated. If one replaces B by \(B/Ann_{B}(N)\) , the hypotheses of the proposition are still satisfied (§ 2, no. 1, example 5); one may therefore suppose that B is a finitely generated A-module and that \(\operatorname{Supp}_{B}(N) = \operatorname{Spec}(B)\) .

Let m be a maximal ideal of Supp \(_A\) (N) ; set \(n = \dim(A_m)\) . By cor. 2 of prop. 3 of no. 1, it suffices to prove that N \(_m\) is a Macaulay A \(_m\) -module of dimension \(n\) . Every maximal ideal of B \(_m\) is of the form nB \(_m\) , where n is a prime ideal of B lying over m (V, § 2, no. 1, lemma 1 and prop. 1). Let p be a minimal prime ideal of Supp \(_B\) The closed subset V(pB \(_n\) ) of Supp \(_{B_n}\) (N \(_n\) ) is then of codimension zero; the B \(_n\) -module N \(_n\) being Macaulay, the cor. of prop. 2, § 2, no. 2 implies the equality dim \(_{B_n}\) (N \(_n\) ) = dim(B \(_n\) /pB \(_n\) ). But p lies above a minimal prime ideal of A, contained in m, so that \(pB_{n}\) lies above a minimal prime ideal of \(A_{m}\) , which is zero since the local ring \(A_{m}\) is regular, hence integral. The canonical map \(A_{m} \to B_{n}/pB_{n}\) is therefore injective, and it follows from VIII, § 2, no. 3, th. 1 that we have \(\dim(B_{n}/pB_{n}) = n\) Proposition 2 of § 2, no. 2 implies that the \(A_{m}\) -module \(N_{m}\) is Macaulay. The proposition then follows from the relations \(\dim_{A_{m}}(N_{m}) = \dim_{B_{m}}(N_{m})\) (VIII, § 2, no. 3, th. 1) and \(\dim_{B_{m}}(N_{m}) \geqslant \dim_{B_{n}}(N_{n}) = n\) .

COROLLARY.--- Let B be a Noetherian integral domain and let A be a regular subring of B. Suppose that B is a finitely generated A-module. For B to be a Macaulay ring, it is necessary and sufficient that the A-module B be projective.

If the ring B is Macaulay, the A-module B is projective by prop. 5.

Conversely, suppose the A-module B is projective. The A-module A is Macaulay (cor. 2 of prop. 4); the A-module B is a direct summand in a finitely generated free A-module, hence is Macaulay (§ 2, no. 1, example 2). One then applies cor. 1 of prop. 8 of § 2, no. 6.

Example.--- Let B be a nonzero integral domain of finite type over a field K. By VIII, § 2, no. 4, cor. 1 of th. 3, there exists a subalgebra A of B isomorphic to a polynomial algebra over K such that B is a finitely generated A-module. The following properties are equivalent:

\begin{enumerate}
\def\labelenumi{(\roman{enumi})}
\item
  the ring B is Macaulay;
\item
  the A-module B is projective;
\item
  the A-module B is free.
\end{enumerate}

